%%%%%%%%%%%%%%%%%%%%%%%%%%%%%%%%%%%%%%%%
% datoteka diploma_predloga.tex
%
% ta verzija sproducira pdf datoteko v pdf/A formatu.
%
% vzorčna datoteka za pisanje diplomskega dela v formatu LaTeX
% na UL Fakulteti za računalništvo in informatiko
%
% Leta 2026 je predlogo posodobil Žiga Emeršič. Predlogo je razdelil na več datotek, uvedel spremenljivke, dodal PDF-A in posodobil navodila (tudi združil določena poglavja in pripravil nova).
% Leta 2021 je predlogo posodobil Franc Solina.
% Leta 2010 je predlogo pripravil Gašper Fijavž.
%
% za upravljanje z literaturo ta vezija uporablja BibLaTeX
%
% svetujemo uporabo Overleaf.com - na tej spletni implementaciji LaTeXa ta vzorec zagotovo pravilno deluje
%
% -------------------------------------------------------------------
% OPOMBA: izvorna vzorčna vsebina predloge (poglavja o LaTeXu, plovkah,
% matematičnem okolju itd.) je bila odstranjena in nadomeščena z vsebino
% diplomske naloge. Vzorec je še vedno dosegljiv v zgodovini repozitorija:
%     git show 156ad49:diploma-predloga.tex
% -------------------------------------------------------------------

% Združljivost s pdfa (sicer nekoliko starejšo različico 1.7)
\DocumentMetadata{pdfstandard = a-2b, pdfversion  = 1.7}

\documentclass[a4paper,12pt,openright]{book}
%\documentclass[a4paper, 12pt, openright, draft]{book}  Nalogo preverite tudi z opcijo draft, ki pokaže, katere vrstice so predolge! Pozor, v draft opciji, se slike ne pokažejo!


%%%%%%%%%%%%%%%%%%%%%%%%%%%%%%%%%%%%%%%%
%  PODATKI O DIPLOMI
%%%%%%%%%%%%%%%%%%%%%%%%%%%%%%%%%%%%%%%%

% Vnesite vaše podatke in podatke mentorjev:
\newcommand{\tauthor}{Tilen Zrnec}
\newcommand{\temail}{tilen.zrnec@student.uni-lj.si}
\newcommand{\tmentor}{izr. prof. dr. Matjaž Kukar}
\newcommand{\tsomentor}{/}

% Odkomentirajte pravi program:
\newcommand{\tstopnja}%
% {VISOKOŠOLSKI STROKOVNI ŠTUDIJSKI PROGRAM\\ PRVE STOPNJE}
{UNIVERZITETNI  ŠTUDIJSKI PROGRAM\\ PRVE STOPNJE\\ RAČUNALNIŠTVO IN INFORMATIKA}
% {INTERDISCIPLINARNI UNIVERZITETNI\\ ŠTUDIJSKI PROGRAM PRVE STOPNJE\\ MULTIMEDIJA}
% {INTERDISCIPLINARNI UNIVERZITETNI\\ ŠTUDIJSKI PROGRAM PRVE STOPNJE\\ UPRAVNA INFORMATIKA}
% {INTERDISCIPLINARNI UNIVERZITETNI\\ ŠTUDIJSKI PROGRAM PRVE STOPNJE\\ RAČUNALNIŠTVO IN MATEMATIKA}

% Odkomentirajte pravo institucijo:
\newcommand{\tinst}%
{Fakulteta za računalništvo in informatiko}
% {Fakulteta za elektrotehniko}
% {Fakulteta za upravo}
% {Fakulteta za matematiko in fiziko}

% Napišite pravo vrsto naloge:
\newcommand{\tvrsta}{Diplomska naloga na univerzitetnem študijskem programu prve stopnje Računalništvo in informatika}

% Opis iz STUDISa:
% TODO: prepiši besedilo teme diplomskega dela iz sistema STUDIS (vnesel ga je mentor).
\newcommand{\tdescription}
{Besedilo teme diplomskega dela študent prepiše iz študijskega informacijskega sistema, kamor ga je vnesel mentor. V nekaj stavkih bo opisal, kaj pričakuje
od kandidatovega diplomskega dela. Kaj so cilji, kakšne metode naj uporabi,
morda bo zapisal tudi ključno literaturo.}
\newcommand{\tdescriptionEn}
{The student shall transcribe the text of the diploma thesis topic from the student information system, where it was entered by the mentor. In a few sentences, they will describe what is expected of the candidate's thesis. What the objectives are, what methods should be used, and perhaps they will also list the key literature.}

% Naslov:
\newcommand{\ttitle}{Tabelarični temeljni modeli}
\newcommand{\ttitleEn}{Tabular Foundation Models}

% Ključne besede:
\newcommand{\tkeywords}{tabelarični temeljni modeli, strojno učenje, tabelarični podatki, ansambli odločitvenih dreves, učenje iz konteksta, OpenML-CC18}
\newcommand{\tkeywordsEn}{tabular foundation models, machine learning, tabular data, decision tree ensembles, in-context learning, OpenML-CC18}

% Zahvala (opcijsko):
% TODO: napiši zahvalo ali celoten ukaz zakomentiraj.
\newcommand{\tzahvala}
{Na tem mestu zapišite, komu se zahvaljujete za pomoč pri izdelavi diplomske naloge oziroma pri vašem študiju nasploh. Pazite, da ne boste koga pozabili. Utegnil vam bo zameriti. Temu se da izogniti tako, da celotno zahvalo izpustite.}

% Posvetilo (opcijsko):
\newcommand{\tposvetilo}
{Svojemu dragemu računalniku.}

%%%%%%%%%%%%%%%%%%%%%%%%%%%%%%%%%%%%%%%%
%  VKLJUČITEV NUJNIH ELEMENTOV 1/2
%%%%%%%%%%%%%%%%%%%%%%%%%%%%%%%%%%%%%%%%

%%%%%%%%%%%%%%%%%%%%%%%%%%%%%%%%%%%%%%%%
% datoteka preambula.tex se uporablja v diploma-FRI-vzorec.tex
%%%%%%%%%%%%%%%%%%%%%%%%%%%%%%%%%%%%%%%%

\usepackage[T1]{fontenc}      % omogoča normalno kopiranje slovenskih črk
\usepackage[utf8]{inputenc}   % omogoča uporabo slovenskih črk kodiranih v formatu UTF-8
\usepackage[slovene,american]{babel}    % naloži, med drugim, slovenske naslove poglavij in delilne vzorce
\usepackage{microtype}         % boljša postavitev
\usepackage[pdftex]{graphicx}  % omogoča vlaganje slik različnih formatov
\usepackage{fancyhdr}          % poskrbi, na primer, za glave strani
\usepackage{amssymb}           % dodatni matematični simboli
\usepackage{amsmath}           % eqref, npr.
\usepackage{amsfonts}
\usepackage[hyphens]{url}
\usepackage{color}
\usepackage{soul}
\usepackage{subcaption}
\usepackage{algorithm}
\usepackage{algpseudocode}
\usepackage[
  backend=biber,
  style=numeric,
  sorting=nty,
]{biblatex}

\addbibresource{literatura.bib}        % pot do literature
\graphicspath{ {./podporno/}{./slike/} } % poti do slik

%%%%%%%%%%%%%%%%%%%%%%%%%%%%%%%%%%%%%%%%
%  enote (decimalna vejica)
%%%%%%%%%%%%%%%%%%%%%%%%%%%%%%%%%%%%%%%%

\usepackage{siunitx}
\sisetup{
  output-decimal-marker = {,}
}

%%%%%%%%%%%%%%%%%%%%%%%%%%%%%%%%%%%%%%%%
%  URLs
%%%%%%%%%%%%%%%%%%%%%%%%%%%%%%%%%%%%%%%%

% najprej podtaknemo opcijo za deljenje URL-jev pri vezajih
\PassOptionsToPackage{hyphens}{url}

% nato naložimo hyperref z vsemi tvojimi nastavitvami
\usepackage[
    pdftex, 
    unicode,                % Podpora za slovenske znake v PDF kazalu/metapodatkih
    colorlinks=true,
    citecolor=black, 
    filecolor=black,
    linkcolor=black, 
    urlcolor=black,
    pdfproducer={LaTeX}, 
    pdfcreator={LaTeX}
]{hyperref}


% \newcommand{\ttitle}{Vzorec diplomskega dela}
% \newcommand{\ttitleEn}{Diploma thesis template}
% \newcommand{\tsubject}{\ttitle}
% \newcommand{\tsubjectEn}{\ttitleEn}
% \newcommand{\tauthor}{Matjaž Kralj}
% \newcommand{\tkeywords}{računalnik, računalnik, računalnik}
% \newcommand{\tkeywordsEn}{computer, computer, computer}

%%%%%%%%%%%%%%%%%%%%%%%%%%%%%%%%%%%%%%%%
%  HYPERREF SETUP
%%%%%%%%%%%%%%%%%%%%%%%%%%%%%%%%%%%%%%%%
\hypersetup{pdftitle={\ttitle}}
\hypersetup{pdfsubject=\ttitleEn}
\hypersetup{pdfauthor={\tauthor}}
\hypersetup{pdfkeywords=\tkeywordsEn}

%%%%%%%%%%%%%%%%%%%%%%%%%%%%%%%%%%%%%%%%
% postavitev strani
%%%%%%%%%%%%%%%%%%%%%%%%%%%%%%%%%%%%%%%%

\addtolength{\marginparwidth}{-20pt} % robovi za tisk
\addtolength{\oddsidemargin}{40pt}
\addtolength{\evensidemargin}{-40pt}

\renewcommand{\baselinestretch}{1.3} % ustrezen razmik med vrsticami
\setlength{\headheight}{15pt}        % potreben prostor na vrhu
\renewcommand{\chaptermark}[1]%
{\markboth{\MakeUppercase{\thechapter.\ #1}}{}} \renewcommand{\sectionmark}[1]%
{\markright{\MakeUppercase{\thesection.\ #1}}} \renewcommand{\headrulewidth}{0.5pt} \renewcommand{\footrulewidth}{0pt}
\fancyhf{}
\fancyhead[LE,RO]{\sl \thepage}
%\fancyhead[LO]{\sl \rightmark} \fancyhead[RE]{\sl \leftmark}
\fancyhead[RE]{\sc \tauthor}              % dodal Solina
\fancyhead[LO]{\sc Diplomska naloga}     % dodal Solina

\newcommand{\BibLaTeX}{{\sc Bib}\LaTeX}
\newcommand{\BibTeX}{{\sc Bib}\TeX}

%%%%%%%%%%%%%%%%%%%%%%%%%%%%%%%%%%%%%%%%
% naslovi
%%%%%%%%%%%%%%%%%%%%%%%%%%%%%%%%%%%%%%%%

\newcommand{\autfont}{\Large}
\newcommand{\titfont}{\LARGE\bf}
\newcommand{\clearemptydoublepage}{\newpage{\pagestyle{empty}\cleardoublepage}}
\setcounter{tocdepth}{1}        % globina kazala

%%%%%%%%%%%%%%%%%%%%%%%%%%%%%%%%%%%%%%%%
% konstrukti
%%%%%%%%%%%%%%%%%%%%%%%%%%%%%%%%%%%%%%%%
\newtheorem{izrek}{Izrek}[chapter]
\newtheorem{trditev}{Trditev}[izrek]
\newenvironment{dokaz}{\emph{Dokaz.}\ }{\hspace{\fill}{$\Box$}}

%%%%%%%%%%%%%%%%%%%%%%%%%%%%%%%%%%%%%%%%
% define medatata
%%%%%%%%%%%%%%%%%%%%%%%%%%%%%%%%%%%%%%%%
\def\Title{\ttitle}
\def\Author{\tauthor, \temail}
\def\Subject{\ttitleEn}
\def\Keywords{\tkeywordsEn}

%%%%%%%%%%%%%%%%%%%%%%%%%%%%%%%%%%%%%%%%
% \convertDate converts D:20080419103507+02'00' to 2008-04-19T10:35:07+02:00
%%%%%%%%%%%%%%%%%%%%%%%%%%%%%%%%%%%%%%%%
\def\convertDate{%
  \getYear
}

{\catcode`\D=12
  \gdef\getYear D:#1#2#3#4{\edef\xYear{#1#2#3#4}\getMonth}
}
\def\getMonth#1#2{\edef\xMonth{#1#2}\getDay}
\def\getDay#1#2{\edef\xDay{#1#2}\getHour}
\def\getHour#1#2{\edef\xHour{#1#2}\getMin}
\def\getMin#1#2{\edef\xMin{#1#2}\getSec}
\def\getSec#1#2{\edef\xSec{#1#2}\getTZh}
\def\getTZh +#1#2{\edef\xTZh{#1#2}\getTZm}
\def\getTZm '#1#2'{%
  \edef\xTZm{#1#2}%
  \edef\convDate{\xYear-\xMonth-\xDay T\xHour:\xMin:\xSec+\xTZh:\xTZm}%
}

%\expandafter\convertDate\pdfcreationdate

%%%%%%%%%%%%%%%%%%%%%%%%%%%%%%%%%%%%%%%%
% get pdftex version string
%%%%%%%%%%%%%%%%%%%%%%%%%%%%%%%%%%%%%%%%
\newcount\countA
\countA=\pdftexversion
\advance \countA by -100
\def\pdftexVersionStr{pdfTeX-1.\the\countA.\pdftexrevision}

%%%%%%%%%%%%%%%%%%%%%%%%%%%%%%%%%%%%%%%%
% XMP data
%%%%%%%%%%%%%%%%%%%%%%%%%%%%%%%%%%%%%%%%
\usepackage{xmpincl}
%\includexmp{pdfa-1b}

%%%%%%%%%%%%%%%%%%%%%%%%%%%%%%%%%%%%%%%%
% pdfInfo
%%%%%%%%%%%%%%%%%%%%%%%%%%%%%%%%%%%%%%%%
\pdfinfo{%
  /Title    (\ttitle)
  /Author   (\tauthor)
  /Subject  (\ttitleEn)
  /Keywords (\tkeywordsEn)
  /ModDate  (\pdfcreationdate)
  /Trapped  /False
}

%%%%%%%%%%%%%%%%%%%%%%%%%%%%%%%%%%%%%%%%
% vključevanje programske kode
%%%%%%%%%%%%%%%%%%%%%%%%%%%%%%%%%%%%%%%%
\usepackage[hypcap=false]{caption}
%\usepackage{minted}
\definecolor{svetlo-siva}{gray}{0.96}
%\usemintedstyle{friendly} % lahko uporabite kakšno drugo temo
%\setminted{
%    fontsize=\footnotesize,
%    baselinestretch=1.1,
%    linenos,            % številčenje vrstic
%    frame=lines,        % črte pred in za
%    framesep=2mm,
%    bgcolor=svetlo-siva, 
%    breaklines=true     % lomljenje dolgih vrstic
%}

% izsek za slovenščino
%\addto\captionsslovene{%
%  \renewcommand{\listingscaption}{Izsek kode}%
%}

% izsek za angleščino
%\addto\captionsenglish{%
%  \renewcommand{\listingscaption}{Code Snippet}%
%}

%%%%%%%%%%%%%%%%%%%%%%%%%%%%%%%%%%%%%%%%
% znaki za copyright stran
%%%%%%%%%%%%%%%%%%%%%%%%%%%%%%%%%%%%%%%%

\newcommand{\CcImageCc}[1]{%
  \includegraphics[scale=#1]{cc_cc_30.pdf}%
}
\newcommand{\CcImageBy}[1]{%
  \includegraphics[scale=#1]{cc_by_30.pdf}%
}
\newcommand{\CcImageSa}[1]{%
  \includegraphics[scale=#1]{cc_sa_30.pdf}%
}

%%%%%%%%%%%%%%%%%%%%%%%%%%%%%%%%%%%%%%%%
\usepackage[autostyle=true]{csquotes} % za narekovaje

% Če želite uporabljati klasične slovenske „narekovaje“ (so enaki nemškim):
\DeclareQuoteAlias{german}{slovene}
% Če pa želite imeti koničaste tipkarske »narekovaje« pa zgornjo vrstico odstranite.


%%%%%%%%%%%%%%%%%%%%%%%%%%%%%%%%%%%%%%%%
%  POMOŽNI UKAZI ZA ČAS PISANJA
%%%%%%%%%%%%%%%%%%%%%%%%%%%%%%%%%%%%%%%%
% Vidna oznaka za nedokončane dele. Pred oddajo morajo iz naloge izginiti
% vsi klici \todo -- poišči jih z:  grep -n "todo{" diploma-predloga.tex
\newcommand{\todo}[1]{%
  \medskip\noindent\fbox{\parbox{0.95\linewidth}{\textbf{TODO:} \textit{#1}}}\medskip%
}

\begin{document}

%%%%%%%%%%%%%%%%%%%%%%%%%%%%%%%%%%%%%%%%
% POVZETEK
%%%%%%%%%%%%%%%%%%%%%%%%%%%%%%%%%%%%%%%%
% Povzetek in abstract pišemo čisto na koncu, ko so znani vsi rezultati.
% Obsegati morata 100-200 besed in vsebovati: (1) stanje področja in problem,
% (2) naš pristop, (3) najpomembnejši rezultat.

\newcommand{\tpovzetek}
{TODO: povzetek napiši na koncu, ko bodo znani vsi rezultati. Obsegati mora med 100 in 200 besed ter zajeti stanje področja in obravnavani problem, naš pristop k njegovemu reševanju ter najpomembnejši rezultat naloge.}

\newcommand{\tabstract}
{TODO: write the abstract last, once all results are known. It must contain between 100 and 200 words and cover the state of the art and the problem addressed, our approach to solving it, and the most important result of the thesis.}

%%%%%%%%%%%%%%%%%%%%%%%%%%%%%%%%%%%%%%%%
%  VKLJUČITEV NUJNIH ELEMENTOV 2/2
%%%%%%%%%%%%%%%%%%%%%%%%%%%%%%%%%%%%%%%%

%%%%%%%%%%%%%%%%%%%%%%%%%%%%%%%%%%%%%%%%
% datoteka naslovnestrani.tex se uporablja v diploma-FRI-vzorec.tex
%%%%%%%%%%%%%%%%%%%%%%%%%%%%%%%%%%%%%%%%

\selectlanguage{slovene}
\frontmatter
\setcounter{page}{1} %
\renewcommand{\thepage}{}       % preprečimo težave s številkami strani v kazalu

%%%%%%%%%%%%%%%%%%%%%%%%%%%%%%%%%%%%%%%%
%naslovnica
\thispagestyle{empty}%
\begin{center}
  {\large\sc Univerza v Ljubljani\\%
    \tinst\\%
  }
  \vskip 10em%
  {\autfont \tauthor\par}%
  {\titfont \ttitle \par}%
  {\vskip 3em \textsc{DIPLOMSKO DELO\\[5mm]         % dodal Solina za ostale študijske programe
         \tstopnja}\par}%
  \vfill\null%
  % izberite pravi habilitacijski naziv mentorja!
  {\large \textsc{Mentor}: \tmentor\par}%
  {\large \textsc{Somentor}:  \tsomentor\par}%
  {\vskip 2em \large Ljubljana, \the\year \par}%
\end{center}
% prazna stran
%\clearemptydoublepage
% izjava o licencah itd. se izpiše na hrbtni strani naslovnice

%%%%%%%%%%%%%%%%%%%%%%%%%%%%%%%%%%%%%%%%
%copyright stran
%%%%%%%%%%%%%%%%%%%%%%%%%%%%%%%%%%%%%%%%
\newpage
\thispagestyle{empty}

\vspace*{5cm}
{\small \noindent
  To delo je ponujeno pod licenco \textit{Creative Commons Priznanje avtorstva-Deljenje pod enakimi pogoji 2.5 Slovenija} (ali novej\v so razli\v cico).
  To pomeni, da se tako besedilo, slike, grafi in druge sestavine dela kot tudi rezultati diplomskega dela lahko prosto distribuirajo,
  reproducirajo, uporabljajo, priobčujejo javnosti in predelujejo, pod pogojem, da se jasno in vidno navede avtorja in naslov tega
  dela in da se v primeru spremembe, preoblikovanja ali uporabe tega dela v svojem delu, lahko distribuira predelava le pod
  licenco, ki je enaka tej.
  Podrobnosti licence so dostopne na spletni strani \href{http://creativecommons.si}{creativecommons.si} ali na Inštitutu za
  intelektualno lastnino, Streliška 1, 1000 Ljubljana.

  \vspace*{1cm}
  \begin{center}% 0.66 / 0.89 = 0.741573033707865
    \CcImageCc{0.741573033707865}\hspace*{1ex}\CcImageBy{1}\hspace*{1ex}\CcImageSa{1}%
  \end{center}
}

\vspace*{1cm}
{\small \noindent
  Izvorna koda diplomskega dela, njeni rezultati in v ta namen razvita programska oprema je ponujena pod licenco GNU General Public License,
  različica 3 (ali novejša). To pomeni, da se lahko prosto distribuira in/ali predeluje pod njenimi pogoji.
  Podrobnosti licence so dostopne na spletni strani \url{http://www.gnu.org/licenses/}.
}

\vfill
\begin{center}
  \ \\ \vfill
  {\em
  Besedilo je oblikovano z urejevalnikom besedil \LaTeX.}
\end{center}

% prazna stran
\clearemptydoublepage

%%%%%%%%%%%%%%%%%%%%%%%%%%%%%%%%%%%%%%%%
% stran 3 med uvodnimi listi
\thispagestyle{empty}
\
\vfill

\bigskip
\noindent\textbf{Kandidat:} \tauthor\\
\noindent\textbf{Naslov:} \ttitle\\
% vstavite ustrezen naziv študijskega programa!
\noindent\textbf{Vrsta naloge:} \tvrsta \\
% izberite pravi habilitacijski naziv mentorja!
\noindent\textbf{Mentor:} \tmentor \\
\noindent\textbf{Somentor:} \tsomentor

\bigskip
\noindent\textbf{Opis:}\\
\tdescription

\bigskip
\noindent\textbf{Title:} \ttitleEn

\bigskip
\noindent\textbf{Description:}\\
\tdescriptionEn

\vfill

\vspace{2cm}

% prazna stran
\clearemptydoublepage

% zahvala
\thispagestyle{empty}\mbox{}\vfill\null\it%
\noindent
\tzahvala
\rm\normalfont

% prazna stran
\clearemptydoublepage

%%%%%%%%%%%%%%%%%%%%%%%%%%%%%%%%%%%%%%%%
% posvetilo, če sama zahvala ne zadošča :-)
\thispagestyle{empty}\mbox{}{\vskip0.20\textheight}\mbox{}\hfill
\begin{minipage}{0.55\textwidth}%
  \tposvetilo
  \normalfont
\end{minipage}

% prazna stran
\clearemptydoublepage

%%%%%%%%%%%%%%%%%%%%%%%%%%%%%%%%%%%%%%%%
% kazalo
\pagestyle{empty}
\def\thepage{}% preprečimo težave s številkami strani v kazalu
\tableofcontents{}

% prazna stran
\clearemptydoublepage

%%%%%%%%%%%%%%%%%%%%%%%%%%%%%%%%%%%%%%%%
% povzetek
\addcontentsline{toc}{chapter}{Povzetek}
\chapter*{Povzetek}

\noindent\textbf{Naslov:} \ttitle
\bigskip

\noindent\textbf{Avtor:} \tauthor
\bigskip

\noindent \tpovzetek

\bigskip

\noindent\textbf{Ključne besede:} \tkeywords.
% prazna stran
\clearemptydoublepage

%%%%%%%%%%%%%%%%%%%%%%%%%%%%%%%%%%%%%%%%
% abstract
\selectlanguage{english}
\addcontentsline{toc}{chapter}{Abstract}
\chapter*{Abstract}

\noindent\textbf{Title:} \ttitleEn
\bigskip

\noindent\textbf{Author:} \tauthor
\bigskip

\noindent \tabstract
\bigskip

\noindent\textbf{Keywords:} \tkeywordsEn.
\selectlanguage{slovene}
% prazna stran
\clearemptydoublepage
\mainmatter
\setcounter{page}{1}
\pagestyle{fancy}

%%%%%%%%%%%%%%%%%%%%%%%%%%%%%%%%%%%%%%%%
% GLAVNA VSEBINA
%%%%%%%%%%%%%%%%%%%%%%%%%%%%%%%%%%%%%%%%

%=========================================================================
\chapter{Uvod}
\label{ch:uvod}
%=========================================================================
% Vsebina poglavja (piši ga med zadnjimi, ko so rezultati znani):
%  - Zakaj so tabelarični podatki pomembni in zakaj so dolgo veljali za
%    področje, kjer globoko učenje zaostaja za ansambli odločitvenih dreves.
%  - Kaj so temeljni modeli in kaj obljubljajo na tabelaričnih podatkih.
%  - Motivacija: ali napovedni model, ki se ga sploh ne uči na ciljni
%    podatkovni množici, prekaša uveljavljene ansamble?
%  - Cilji naloge in raziskovalna vprašanja.
%  - Prispevki naloge.
%  - Zemljevid naloge: kratek opis vsebine po poglavjih (z ukazom \ref,
%    številk poglavij nikoli ne piši ročno).

\todo{Napiši uvod: motivacija, cilji, raziskovalna vprašanja, prispevki in zemljevid naloge.}

%=========================================================================
\chapter{Sorodni pristopi}
\label{ch:sorodni}
%=========================================================================
% Vsebina poglavja:
%  - Klasični pristopi k tabelaričnim podatkom: ansambli odločitvenih dreves
%    kot dolgoletno privzeto orodje.
%  - Globoko učenje na tabelaričnih podatkih in ugotovitve primerjalnih
%    študij, da drevesni modeli pogosto ostajajo boljši.
%  - Temeljni modeli in učenje iz konteksta; prenos ideje na tabelarične
%    podatke (TabPFN, TabICL in sorodni modeli).
%  - Zbirke za primerjalno vrednotenje tabelaričnih modelov (OpenML-CC18,
%    TabArena) in metodološke kritike obstoječih primerjav.
%  - Umestitev: v čem se naša primerjava razlikuje od objavljenih.

\todo{Napiši pregled sorodnih pristopov in umesti nalogo v obstoječi kontekst.}

%=========================================================================
\chapter{Metode in testna metodologija}
\label{ch:metode}
%=========================================================================
% Opis podatkovnih množic ni v tem poglavju, ampak v poglavju \ref{ch:podatki},
% ki prevzema vlogo prvotno načrtovanega podpoglavja 3.1 (Podatkovna nabora).
% To poglavje opisuje algoritme in način testiranja.

\section{Algoritmi}
\label{sec:algoritmi}

\todo{Six subsections, roughly a page each: what it is, why it's in the comparison, version pinned, the defaults that matter. This is your longest section and probably your best use of today.}

\section{Protokol vrednotenja}
\label{sec:protokol}

\todo{5-fold stratified CV, identical splits, \texttt{random\_state=42}, ROC-AUC binary vs.\ macro-OVR. This is the spine everything else hangs off.}

\section{Politika predobdelave}
\label{sec:predobdelava}

\todo{Your experimental variable. Say explicitly that it's a deliberate design choice, not an oversight, and that it's what Medic3 will finally stress.}

\section{Izvedbeno okolje}
\label{sec:okolje}

\todo{Arnes, H100, SLURM job arrays, micromamba, per-fit checkpointing with atomic \texttt{.partial} files, fail-soft error handling. Fast, descriptive writing. Good warm-up if the algorithms section feels heavy.}

\section{Ponovljivost}
\label{sec:ponovljivost}

\todo{Repo, seeds, pinned versions, PROVENANCE.}

%=========================================================================
\chapter{Materiali}
\label{ch:podatki}
%=========================================================================
% Vsebina poglavja:
%  - Zbirka OpenML-CC18: kako je nastala, kriteriji za uvrstitev množic,
%    zakaj je primerna za to primerjavo (uporabljata jo tudi izvirna članka
%    o modelih TabPFN in TabICL).
%  - Opisna statistika vseh 72 podatkovnih množic: število primerov,
%    število atributov, delež kategoričnih atributov, število razredov,
%    delež manjkajočih vrednosti, uravnoteženost razredov.
%  - Tabela z lastnostmi množic (lahko v prilogo, v besedilu pa povzetek).
%  - Slika: porazdelitev velikosti množic (log-log razsevni diagram
%    število primerov proti številu atributov).
%  - Kje se množice približajo mejam modelov (npr. največje množice).
%
% To poglavje prevzema vlogo prvotno načrtovanega podpoglavja 3.1
% (Podatkovna nabora).

\todo{CC18 half is writable now (inclusion criteria, 72 datasets, size ranges). Leave the Medic3 subsection as a stub until after the Zoom.}

\section{OpenML-CC18}
\label{sec:cc18}

\todo{Kriteriji za uvrstitev podatkovnih množic v zbirko, 72 množic in razponi velikosti: število primerov, število atributov, delež kategoričnih atributov, število razredov in delež manjkajočih vrednosti. Tabela z opisno statistiko in slika porazdelitve velikosti.}

\section{Bolnišnična podatkovna množica Medic3}
\label{sec:medic3}

\todo{Ostane zgolj okostje do pogovora prek Zooma.}

%=========================================================================
\chapter{Rezultati}
\label{ch:rezultati}
%=========================================================================
% Vsebina poglavja (samo predstavitev meritev, razlaga gre v poglavje 6):
%  - Skupna primerjava: povprečni ROC-AUC in povprečni rang po algoritmih.
%  - Tabela z rezultati po posameznih podatkovnih množicah.
%  - Statistična značilnost razlik med pari algoritmov.
%  - Časovna učinkovitost: čas učenja in čas sklepanja.
%  - Odvisnost uspešnosti od velikosti podatkovne množice in od deleža
%    kategoričnih atributov.
%  - Neuspela učenja in njihovi vzroki.

\todo{Predstavi rezultate: skupna primerjava, rezultati po množicah, statistična značilnost, časi in neuspela učenja.}

%=========================================================================
\chapter{Diskusija in zaključek}
\label{ch:zakljucek}
%=========================================================================
% Vsebina poglavja:
%  - Razlaga rezultatov in odgovori na raziskovalna vprašanja iz uvoda.
%  - Primerjava naših ugotovitev z objavljenimi rezultati.
%  - Praktična priporočila: kdaj se splača uporabiti temeljni model.
%  - Omejitve dela (privzeti hiperparametri, ena metrika, omejitve
%    strojne opreme, zgolj klasifikacija).
%  - Nadaljnje delo: optimizacija hiperparametrov, bolnišnična podatkovna
%    množica Medic3, dodatne zbirke za primerjalno vrednotenje.

\todo{Napiši diskusijo in zaključek: razlaga rezultatov, omejitve in nadaljnje delo.}

%%%%%%%%%%%%%%%%%%%%%%%%%%%%%%%%%%%%%%%%
%  LITERATURA
%%%%%%%%%%%%%%%%%%%%%%%%%%%%%%%%%%%%%%%%

%\cleardoublepage
\printbibliography[heading=bibintoc]

\end{document}
